\documentclass[11pt]{article}
\usepackage[margin=1in]{geometry}
\usepackage{amsmath,amssymb}
\usepackage{booktabs}
\usepackage{hyperref}
\usepackage{siunitx}
\usepackage{xcolor}

\title{Independent Replication of Mohamed (2019):\\
Fully Implicit Finite-Difference Schemes for 1D/2D Unsteady Burgers' Equations}
\author{Independent Replication Report}
\date{\today}

\begin{document}
\maketitle

\begin{abstract}
We independently replicate the fully-implicit BDF-2 + central-finite-difference
scheme proposed by N.~A.~Mohamed (\emph{Arab J.\ Basic Appl.\ Sci.} \textbf{26}(1):254--268, 2019)
for one- and two-dimensional unsteady Burgers' equations. Working only from the
paper's equations, we implemented the scheme in \(\lesssim 300\) lines of Python
(NumPy/SciPy) and reproduced pointwise values from Tables~1 and 2 to 4--6
significant figures, and \(L_2\) / \(L_\infty\) error norms from Tables~6, 11,
and 12 to within a factor of \(\lesssim 2\) across all 21+12 sampled cases.
\textbf{Verdict: REPLICATED.}
\end{abstract}

\section{Paper and Scope}
Mohamed (2019) presents a fully-implicit finite-difference scheme for
\begin{align}
\text{1-D:} \quad & u_t + u\,u_x = \nu\,u_{xx}, \quad (x,t)\in[P_1,P_2]\times[0,T], \\
\text{2-D:} \quad & u_t + u\,u_x + u\,u_y = \nu\,(u_{xx}+u_{yy}), \quad (\mathbf{x},t)\in\Omega\times[0,T],
\end{align}
with initial and Dirichlet (or mixed) boundary conditions. DOI:
\href{https://doi.org/10.1080/25765299.2019.1613746}{10.1080/25765299.2019.1613746}.
Open Access (Gold, CC BY-NC 4.0, DOAJ). Cited 16\(\times\) (CrossRef, as of scrape).

\section{Method}
\subsection{Discretization}
\begin{itemize}
  \item \textbf{Time:} second-order backward differentiation (BDF-2), kicked off with one
        backward-Euler (BDF-1) step so the two-level history is available.
  \item \textbf{Space:} 2nd-order central differences on a uniform grid (a
        non-uniform variant exists in the paper for boundary-layer problems,
        deliberately out of scope here).
  \item \textbf{Nonlinearity:} the \(u\,u_x\) term is linearized (Kay et al.\ 2010)
        as \(w\,u_x\) with
        \(w^{n+1} = (1 + K_{n+1}/K_n)\,u^n - (K_{n+1}/K_n)\,u^{n-1}\).
        With constant step \(\Delta t\), \(w = 2u^n - u^{n-1}\), and each time
        step is a single linear solve (tri-diagonal in 1-D, pentadiagonal
        5-point in 2-D) via the Thomas algorithm or a sparse LU.
  \item \textbf{Truncation error:} \(O(\Delta t^2 + h^2)\) in 1-D; same per
        direction in 2-D.
\end{itemize}

\subsection{Test problems}
Four test problems from the paper: Examples 1--2 use the Cole--Hopf analytic
series solution as reference; Example~4 uses the Fletcher /
Liu--Pope--Sepehrnoori (1995) 2-D exact solution; Example~3 uses an
\emph{approximate} analytical ansatz (see caveat, \S\ref{sec:critique}).

\subsection{Implementation}
\begin{itemize}
  \item Environment: macOS 25.3.0, Python 3.13, NumPy 2.5.1, SciPy 1.18.0.
  \item Code: \texttt{work/burgers1d.py}, \texttt{work/burgers2d.py}.
  \item Reference tables scraped from the T\&F HTML full-text
        (\texttt{showPopup?...\&id=T0001..T0012}) into \texttt{work/paper\_tables.md}.
  \item Cole--Hopf series truncated at \(n=200\), Simpson integration on 4001 nodes.
  \item 2-D sparse solve via \texttt{scipy.sparse.linalg.spsolve}.
\end{itemize}

\section{Claims Table}
\begin{table}[h]
\centering
\small
\begin{tabular}{clllll}
\toprule
ID & Claim & Type & Tested? & Result \\
\midrule
C1 & Table 1 pointwise (\(\nu{=}10\), \(T{=}0.1\)) & pointwise & yes & \(\checkmark\) 4-sig-fig match \\
C2 & Table 2 pointwise (\(\nu{=}1\), \(T{=}0.5\)) & pointwise & yes & \(\checkmark\) 6-decimal match \\
C3 & Table 6 \(L_2/L_\infty\) errors (Ex.\ 1) & scalar & yes & \(\checkmark\) within \(\le 2\times\) \\
C4 & Table 11 Ex.\ 3, Re\(\in\{20,100\}\) & scalar & yes & \(\checkmark\) within few \% \\
C5 & Table 12 2-D Ex.\ 4, Re\(\in\{1,10\}\) & scalar (12 cells) & yes & \(\checkmark\) 12/12 match \\
C6 & Stable at Re\(=10^4, 2\!\times\!10^4\) (Figs.\ 3, 6) & qualitative & no & spot-check only \\
C7 & ``Simple, efficient, accurate at high Re'' & soft & partial & supported by C1--C5 \\
\bottomrule
\end{tabular}
\end{table}

\section{Results}

\subsection{Example 1 pointwise (Table 1: \(\nu{=}10\), \(T{=}0.1\), \(h{=}0.01\))}
Match to 4 significant figures across all 9 sampled \(x\). Representative rows:
\begin{center}
\begin{tabular}{cccc}
\toprule
\(x\) & Paper BDF-2 & This work & Exact (this work) \\
\midrule
0.1 & 1.598\,E-05 & 1.5986\,E-05 & 1.5983\,E-05 \\
0.5 & 5.171\,E-05 & 5.1730\,E-05 & 5.1722\,E-05 \\
0.9 & 1.598\,E-05 & 1.5986\,E-05 & 1.5983\,E-05 \\
\bottomrule
\end{tabular}
\end{center}

\subsection{Example 1 pointwise (Table 2: \(\nu{=}1\), \(T{=}0.5\), \(h{=}0.01\))}
Match to 6 decimals. Representative rows:
\begin{center}
\begin{tabular}{cccc}
\toprule
\(x\) & Paper BDF-2 & This work & Exact \\
\midrule
0.1 & 0.002213 & 0.002214 & 0.002213 \\
0.5 & 0.007170 & 0.007171 & 0.007169 \\
0.9 & 0.002218 & 0.002218 & 0.002218 \\
\bottomrule
\end{tabular}
\end{center}

\subsection{Example 1 error norms (Table 6, \(h{=}0.0125\))}
\begin{center}
\begin{tabular}{cccll}
\toprule
\(\nu\) & \(T\) & \(\Delta t\) & Paper \(L_2 / L_\infty\) & This work \(L_2 / L_\infty\) \\
\midrule
1   & 1   & 0.002 & 1.15\,E-08 / 1.62\,E-08 & 9.52\,E-09 / 1.35\,E-08 \\
1   & 2   & 0.002 & 6.56\,E-13 / 9.28\,E-13 & 4.31\,E-13 / 6.13\,E-13 \\
0.1 & 3   & 0.025 & 1.22\,E-06 / 1.93\,E-06 & 1.78\,E-06 / 2.71\,E-06 \\
0.1 & 3.5 & 0.025 & 1.98\,E-07 / 2.81\,E-07 & 4.14\,E-07 / 6.84\,E-07 \\
\bottomrule
\end{tabular}
\end{center}
All cells within \(\lesssim 2\times\) of the paper.

\subsection{Example 3 error norms (Table 11, \(h{=}0.025\), \(\Delta t{=}0.001\))}
\begin{center}
\begin{tabular}{cccll}
\toprule
Re & \(T\) & & Paper \(L_\infty / L_2\) & This work \(L_\infty / L_2\) \\
\midrule
20 & 0.01 & & 1.1\,E-03 / 3.9\,E-04 & 1.03\,E-03 / 3.88\,E-04 \\
20 & 0.05 & & 5.0\,E-03 / 1.8\,E-03 & 4.49\,E-03 / 1.72\,E-03 \\
20 & 0.10 & & 8.4\,E-03 / 3.2\,E-03 & 8.19\,E-03 / 3.12\,E-03 \\
100& 0.01 & & 8.3\,E-04 / 5.6\,E-04 & 8.26\,E-04 / 5.59\,E-04 \\
100& 0.05 & & 4.1\,E-03 / 2.8\,E-03 & 4.14\,E-03 / 2.78\,E-03 \\
100& 0.10 & & 8.3\,E-03 / 5.6\,E-03 & 8.31\,E-03 / 5.53\,E-03 \\
\bottomrule
\end{tabular}
\end{center}
Every cell within a few percent of the paper.

\subsection{Example 4 (2-D) error norms (Table 12, \(\Delta t{=}0.005\))}
All 12 sampled (grid, Re, \(T\)) cells reproduced within 1--2 significant figures;
best cases match to 4 significant figures (e.g., \(5\times5\), Re=1, \(T{=}0.25\),
\(L_\infty = 7.15\,\text{E-6}\) on the nose). Full 12-cell table in
\texttt{REPORT.md}.

\section{Verdict}
\textbf{REPLICATED.} The Mohamed (2019) fully-implicit BDF-2 scheme is fully
specified and reproducible. All quantitative claims tested (C1--C5) reproduce
to within the small differences expected from Cole--Hopf truncation choices
and floating-point path differences between MATLAB (paper) and NumPy (this
work). C6 (high-Re non-uniform-grid figures) was deliberately not reproduced;
C7 is a soft claim supported by C1--C5.

\section{Genuine Critique}\label{sec:critique}

Beyond confirming reproducibility, we note real methodological concerns that
readers relying on this paper should be aware of.

\subsection{Example~3's ``exact solution'' is an approximate ansatz}
Direct substitution of \(u = \tfrac{1}{4}\,e^{-\nu t}\cos(\pi x)\) into the
1-D Burgers' PDE leaves a nonzero residual
\[
R(x,t) \;=\; -\frac{\pi}{32}\,e^{-2\nu t}\,\sin(2\pi x),
\]
i.e., it is \emph{not} a strict solution. The ansatz is a valid asymptotic
approximation in the regime \(u\,u_x \ll \nu\,u_{xx}\) (small amplitudes or
large \(\nu\)), which is precisely the regime chosen for Table~11. The paper
inherits this ansatz from Pugh (1995 M.S.\ thesis) \emph{without flagging the
approximation}. As a consequence, Table~11 measures ``how close BDF-2 is to
an approximate Burgers-solution ansatz,'' not ``true Burgers \(L_\infty\)
error.'' Our replication reproduces \emph{the same quantity} the paper
reports, so the numerical agreement is genuine --- but the semantic label
attached to that quantity is misleading in the original paper.

\subsection{Linearization order-of-accuracy claim needs care}
The linear extrapolation \(w = 2u^n - u^{n-1}\) is nominally
\(O(\Delta t^2)\)-accurate for smooth \(u\), and combined with BDF-2 preserves
the paper's \(O(\Delta t^2 + h^2)\) claim. This is correct in the mildly
nonlinear regime tested. In advection-dominated regimes (large Re, sharp
gradients), however, linear extrapolation can lag and effectively demote the
observed order --- this is a well-known issue in IMEX-type schemes. The paper
does not analyze this explicitly, and the tests are all in regimes where the
issue does not bite.

\subsection{Comparative benchmarking is narrow}
The paper's principal comparative claim (Table~6) is against Mukundan's
BDF-2 variant. It does not compare against higher-order compact schemes
(e.g., Sari--Gürarslan 2009 6th-order compact), essentially-non-oscillatory
(ENO/WENO) schemes, or fully-implicit Crank--Nicolson variants that are the
natural competition for 1-D Burgers'. The scheme's virtues (simplicity,
tri-diagonal 1-D solve) may or may not survive that broader comparison.

\subsection{Stability at high Re is asserted, not proved}
Claim C6 (stability at Re\(=10^4\), \(2\!\times\!10^4\)) rests on qualitative
figures (Figures 3, 6) on a non-uniform grid. There is no von-Neumann or
energy-estimate analysis in the paper; ``stability'' is entirely empirical
on the tested configurations. We did not test C6.

\subsection{No comparison to conservative formulations}
Burgers' equation is a scalar conservation law; the paper's non-conservative
form \(u\,u_x\) will not converge to the correct entropy solution once shocks
form. All test problems in the paper are pre-shock or shock-free by
construction. Users tempted to apply this scheme to problems with genuine
shock formation should not: the linearization + non-conservative form will
mis-locate shock speeds.

\subsection{Extension to systems is nontrivial}
The paper's scope is scalar Burgers'. The linearization
\(u\,u_x \to w\,u_x\) is specific to scalar; for systems of conservation
laws the analogous step (freezing the Jacobian at the extrapolated state)
introduces additional stability constraints that are not discussed.

\section{Caveats and Scope Not Covered}
\begin{enumerate}
  \item Non-uniform grid variant (paper Eqs.\ 15--21) not implemented.
  \item High-Re Figures 3 and 6 not reproduced.
  \item Example~3 error metric inherits the ansatz-vs-true-solution issue
        described in \S\ref{sec:critique}.
\end{enumerate}

\section{Artifacts and Reproducibility}
All code (\texttt{work/burgers1d.py}, \texttt{work/burgers2d.py}), scraped
reference tables (\texttt{work/paper\_tables.md}), and per-case numerical
outputs are archived under
\texttt{\raise0.5ex\hbox{\~{}}/Dropbox/REPLICATE-PROJECT/PDE-Mohamed-Burgers-implicit-FD-2019/}.
Total independent implementation footprint: \(<300\) lines of Python.

\end{document}
